% GENERATED FILE. Edit canonical lessons, then run npm run generate:books.
% canonical-source-hash: fnv1a64:886cb1ec54d123a7
% canonical-lessons: GE-C513-klavier, GE-C513-geige, GE-C513-floete, GE-C513-schlagzeug, GE-C513-trommeln, GE-R513-review-506-513

\chapter{Klavier und Geige}
\label{ch:ge-klavier-und-geige}
\hlchaptermodality{513}

\begin{chapteropening}
\textbf{By the end of this chapter:} \emph{I can talk about musical instruments.}

Complete the last lesson of chapter 513: I can talk about musical instruments.
\end{chapteropening}

\section[das Klavier]{das Klavier — piano}

\label{lesson:GE-C513-klavier}



\begin{quote}
Before the new one: say the German for ski run, then the German for ski lift.
\end{quote}

\subsection*{You'll want to know: das Klavier}
\textbf{das Klavier} — \textquotedblleft{}piano\textquotedblright{}.

\begin{quote}
Meine Tochter spielt seit drei Jahren Klavier.
\end{quote}

\begin{quote}
\emph{My daughter has been playing the piano for three years.}
\end{quote}

Plural: \emph{die Klaviere}. No article in \emph{Klavier spielen}.

\begin{cousinweb}
\emph{Klavier} was borrowed from French \emph{clavier} \textquotedblleft{}keyboard\textquotedblright{}, which is built on Latin \emph{clāvis} \textquotedblleft{}key\textquotedblright{} (EtymDB).
\end{cousinweb}

\begin{grammarlens}[title={What you've built}]
One more word for telling, explaining and describing at B1.
\end{grammarlens}

\subsection*{Your turn}
\begin{itemize}
\raggedright
  \item \emph{Say it:} \emph{das Klavier}
  \item \emph{Say it:} \emph{das Klavier}, once more
  \item \emph{Say it:} a sentence of your own with \emph{das Klavier}
  \item \emph{Recall:} say the German for ski run, then the German for ski lift, then say \emph{das Klavier} again
\end{itemize}

\begin{tcolorbox}[breakable,colback=teal!4,colframe=teal!35!black,title={Before you move on}]
What does das Klavier mean? (\textquotedblleft{}Piano\textquotedblright{}.) Say it once more.
\end{tcolorbox}
\section[die Geige]{die Geige — violin, fiddle}

\label{lesson:GE-C513-geige}



\begin{quote}
Before the new one: say the German for ski lift, then the German for piano.
\end{quote}

\subsection*{You'll want to know: die Geige}
\textbf{die Geige} — \textquotedblleft{}violin, fiddle\textquotedblright{}.

\begin{quote}
Er spielt Geige in einer kleinen Band.
\end{quote}

\begin{quote}
\emph{He plays the violin in a small band.}
\end{quote}

Plural: \emph{die Geigen}. The formal word is \emph{die Violine}.

\begin{cousinweb}
\emph{Geige} is inherited from Middle High German \emph{gīge} and Old High German \emph{gīga} \textquotedblleft{}fiddle\textquotedblright{} (EtymDB).
\end{cousinweb}

\begin{grammarlens}[title={What you've built}]
One more word for telling, explaining and describing at B1.
\end{grammarlens}

\subsection*{Your turn}
\begin{itemize}
\raggedright
  \item \emph{Say it:} \emph{die Geige}
  \item \emph{Say it:} \emph{die Geige}, once more
  \item \emph{Say it:} a sentence of your own with \emph{die Geige}
  \item \emph{Recall:} say the German for ski lift, then the German for piano, then say \emph{die Geige} again
\end{itemize}

\begin{tcolorbox}[breakable,colback=teal!4,colframe=teal!35!black,title={Before you move on}]
What does die Geige mean? (\textquotedblleft{}Violin, fiddle\textquotedblright{}.) Say it once more.
\end{tcolorbox}
\section[die Flöte]{die Flöte — flute; recorder}

\label{lesson:GE-C513-floete}



\begin{quote}
Before the new one: say the German for piano, then the German for violin.
\end{quote}

\subsection*{You'll want to know: die Flöte}
\textbf{die Flöte} — \textquotedblleft{}flute; recorder\textquotedblright{}.

\begin{quote}
In der Grundschule lernen viele Kinder Flöte.
\end{quote}

\begin{quote}
\emph{Many children learn the recorder at primary school.}
\end{quote}

Plural: \emph{die Flöten}. A recorder is \emph{die Blockflöte}.

\begin{cousinweb}
\emph{Flöte} goes back to Middle High German \emph{floite}, borrowed from Old French \emph{flaüte}, whose own origin is not settled; EtymDB lists the German word without a source.
\end{cousinweb}

\begin{grammarlens}[title={What you've built}]
One more word for telling, explaining and describing at B1.
\end{grammarlens}

\subsection*{Your turn}
\begin{itemize}
\raggedright
  \item \emph{Say it:} \emph{die Flöte}
  \item \emph{Say it:} \emph{die Flöte}, once more
  \item \emph{Say it:} a sentence of your own with \emph{die Flöte}
  \item \emph{Recall:} say the German for piano, then the German for violin, then say \emph{die Flöte} again
\end{itemize}

\begin{tcolorbox}[breakable,colback=teal!4,colframe=teal!35!black,title={Before you move on}]
What does die Flöte mean? (\textquotedblleft{}Flute; recorder\textquotedblright{}.) Say it once more.
\end{tcolorbox}
\section[das Schlagzeug]{das Schlagzeug — drums, drum kit}

\label{lesson:GE-C513-schlagzeug}



\begin{quote}
Before the new one: say the German for violin, then the German for flute.
\end{quote}

\subsection*{You'll want to know: das Schlagzeug}
\textbf{das Schlagzeug} — \textquotedblleft{}drums, drum kit\textquotedblright{}.

\begin{quote}
Der Nachbar übt jeden Abend Schlagzeug.
\end{quote}

\begin{quote}
\emph{The neighbour practises the drums every evening.}
\end{quote}

Plural: \emph{die Schlagzeuge}. One drum is \emph{die Trommel}.

\begin{cousinweb}
\emph{Schlagzeug} is built on \emph{Schlag} \textquotedblleft{}blow\textquotedblright{}, from Proto-Germanic \emph{slagiz}, and \emph{Zeug} \textquotedblleft{}gear, stuff\textquotedblright{}, from Middle High German \emph{ziuc} (EtymDB): the kit you strike.
\end{cousinweb}

\begin{grammarlens}[title={What you've built}]
One more word for telling, explaining and describing at B1.
\end{grammarlens}

\subsection*{Your turn}
\begin{itemize}
\raggedright
  \item \emph{Say it:} \emph{das Schlagzeug}
  \item \emph{Say it:} \emph{das Schlagzeug}, once more
  \item \emph{Say it:} a sentence of your own with \emph{das Schlagzeug}
  \item \emph{Recall:} say the German for violin, then the German for flute, then say \emph{das Schlagzeug} again
\end{itemize}

\begin{tcolorbox}[breakable,colback=teal!4,colframe=teal!35!black,title={Before you move on}]
What does das Schlagzeug mean? (\textquotedblleft{}Drums, drum kit\textquotedblright{}.) Say it once more.
\end{tcolorbox}
\section[trommeln]{trommeln — to drum}

\label{lesson:GE-C513-trommeln}



\begin{quote}
Before the new one: say the German for flute, then the German for drums.
\end{quote}

\subsection*{You'll want to know: trommeln}
\textbf{trommeln} — \textquotedblleft{}to drum\textquotedblright{}.

\begin{quote}
Der Regen trommelt auf das Dach.
\end{quote}

\begin{quote}
\emph{The rain drums on the roof.}
\end{quote}

Regular: \emph{hat getrommelt}. Built on \emph{die Trommel} — drum.

\begin{cousinweb}
\emph{trommeln} is built on \emph{Trommel} \textquotedblleft{}drum\textquotedblright{} (EtymDB lists the noun), usually traced to Middle High German \emph{trumme}, an imitative word for the sound of a drum.
\end{cousinweb}

\begin{grammarlens}[title={What you've built}]
One more word for telling, explaining and describing at B1.
\end{grammarlens}

\subsection*{Your turn}
\begin{itemize}
\raggedright
  \item \emph{Say it:} \emph{trommeln}
  \item \emph{Say it:} \emph{trommeln}, once more
  \item \emph{Say it:} a sentence of your own with \emph{trommeln}
  \item \emph{Recall:} say the German for flute, then the German for drums, then say \emph{trommeln} again
\end{itemize}

\begin{tcolorbox}[breakable,colback=teal!4,colframe=teal!35!black,title={Before you move on}]
What does trommeln mean? (\textquotedblleft{}To drum\textquotedblright{}.) Say it once more.
\end{tcolorbox}
\section[(dialogue)]{Words from chapters 506–513 — a review}

\label{lesson:GE-R513-review-506-513}



\begin{quote}
Say the German for goalkeeper and striker.
\end{quote}

\subsection*{Your turn}
Say these aloud:

\begin{itemize}
\raggedright
  \item goalkeeper, striker, defender, referee, to whistle
  \item half-time, penalty, extra time, tied game, offside
  \item to win, winner, loser, opponent, tournament
  \item competition, championship, trophy, medal, award ceremony
  \item stand, spectator, to cheer on, cheering, to boo
  \item open-air pool, to swim front crawl, to dive, to paddle, canoe
  \item sledge, to toboggan, ice skate, ski run, ski lift
  \item piano, violin, flute, drums, to drum
\end{itemize}

\begin{tcolorbox}[breakable,colback=teal!4,colframe=teal!35!black,title={Before you move on}]
Goalkeeper? (\textbf{der Torwart}.) Striker? (\textbf{der Stürmer}.) Defender? (\textbf{der Verteidiger}.) Referee? (\textbf{der Schiedsrichter}.) To whistle? (\textbf{pfeifen}.) Half-time? (\textbf{die Halbzeit}.) Penalty? (\textbf{der Elfmeter}.) Extra time? (\textbf{die Verlängerung}.) Tied game? (\textbf{das Unentschieden}.) Offside? (\textbf{das Abseits}.) To win? (\textbf{siegen}.) Winner? (\textbf{der Sieger}.) Loser? (\textbf{der Verlierer}.) Opponent? (\textbf{der Gegner}.) Tournament? (\textbf{das Turnier}.) Competition? (\textbf{der Wettkampf}.) Championship? (\textbf{die Meisterschaft}.) Trophy? (\textbf{der Pokal}.) Medal? (\textbf{die Medaille}.) Award ceremony? (\textbf{die Siegerehrung}.) Stand? (\textbf{die Tribüne}.) Spectator? (\textbf{der Zuschauer}.) To cheer on? (\textbf{anfeuern}.) Cheering? (\textbf{der Jubel}.) To boo? (\textbf{buhen}.) Open-air pool? (\textbf{das Freibad}.) To swim front crawl? (\textbf{kraulen}.) To dive? (\textbf{tauchen}.) To paddle? (\textbf{paddeln}.) Canoe? (\textbf{das Kanu}.) Sledge? (\textbf{der Schlitten}.) To toboggan? (\textbf{rodeln}.) Ice skate? (\textbf{der Schlittschuh}.) Ski run? (\textbf{die Piste}.) Ski lift? (\textbf{der Skilift}.) Piano? (\textbf{das Klavier}.) Violin? (\textbf{die Geige}.) Flute? (\textbf{die Flöte}.) Drums? (\textbf{das Schlagzeug}.) To drum? (\textbf{trommeln}.)
\end{tcolorbox}
